% Copyright (c) 2026 Eugene Joseph M. Ragasa.
% SPDX-License-Identifier: Apache-2.0
% Author: Eugene Joseph M. Ragasa
% Formal proof completed: 2026-10-28
\documentclass[11pt]{article}
\usepackage[T1]{fontenc}\usepackage{amsmath,amssymb}\usepackage[margin=1in]{geometry}\usepackage[hidelinks]{hyperref}
\title{Proof Dossier: Finite-Maximum Fourier Truncation Tail Bound}\author{Eugene Joseph M. Ragasa}\date{Formal proof completed: 2026-10-28}
\begin{document}\raggedright\maketitle
\section{Identity and status}
ID \texttt{BG-FOURIER-THM-075}; status \texttt{proved}. Lean declaration:
\par\noindent{\tiny\path{ProjectKoios.BandGauge.Fourier.finiteMaximumNorm_truncationResidual_le_omitted}}.
\section{Statement and proof}
For coefficients $c_\chi$ in a complex normed space and retained support $S$, the finite maximum ambient norm of the deliberate-truncation residual obeys
\[
 \max_k\left\|\sum_{\chi\notin S}\chi(k)c_\chi\right\|
 \leq \sum_{\chi\notin S}\|c_\chi\|.
\]
The omitted coefficient one-norm tail is nonnegative. The finite function-space norm is bounded by a nonnegative scalar exactly when every coordinate norm is bounded by that scalar. Applying \texttt{BG-FOURIER-THM-074} at each momentum therefore yields the maximum bound.
\section{Boundary and audit}
This is an exact finite supremum lift in the ambient norm. It assumes no monotonicity and does not add separately aggregated residual channels. It supplies no entrywise matrix-Frobenius maximum theorem, runtime mesh, binary64 interpretation, physical-range meaning, or scientific-adequacy decision. Axioms are \texttt{propext}, \texttt{Classical.choice}, and \texttt{Quot.sound}; no admissions. Lean \texttt{4.35.0}, Mathlib {\scriptsize\path{56a005855ea17e0b0958f73ec801ae140b0a1679}}; AI-assisted, kernel-checked, no independent expert review recorded.
\end{document}
